\documentclass[11pt,a4paper]{article}
\usepackage[T1]{fontenc}
\usepackage[utf8]{inputenc}
\usepackage[english]{babel}
\usepackage{amsmath,amssymb}
\usepackage[margin=2cm]{geometry}
\setlength{\parindent}{0pt}
\setlength{\emergencystretch}{3em}
\allowdisplaybreaks

\title{About \texttt{false\_critical\_points.tex}}
\author{}
\date{}
\begin{document}
\maketitle
This note explains what the test data in \texttt{false\_critical\_points.tex} is used for. The file is read by \texttt{test\_multivariable\_calculus.c}.

\begin{enumerate}
\item \textbf{Purpose.} The file holds 73 false statements on critical points. A proof of one of them would be the worst error of the engine: a saddle point certified as a minimum. The first twelve are results printed wrong in the collections of solved exercises; the others are true statements of \texttt{critical\_points.tex} changed by little.

\item \textbf{The printed errors.} Item 1: Roma Tre 1(b), the determinant of the Hessian matrix at $(\frac{1}{8}, \frac{1}{4})$ is $-40$, not $-22$. Item 2: Habala 1, at $(-5, 0)$ it is $1080$, not $780$. Item 3: Esparza 2(b), at the origin it is $-36$, not $-4$. Item 4: Eleuteri 5.1.19(5), $(\frac{1}{3}, -\frac{1}{3})$ is a saddle point, not a minimum. Item 5: Perfetti 8.7, $(\frac{1}{\sqrt{2}}, -1)$ is a maximum, not a saddle point. Item 6: Eleuteri 5.1.11, with the printed Hessian matrix $(\frac{1}{2}, \frac{1}{3})$ would be a saddle point; it is a maximum. Item 7: Bosco, $(\frac{2}{3}, -\frac{8}{9})$ ``looks like a minimum''; it is a saddle point. Items 8 and 9: Habala 2, the third leading minors are $-18$ and $18$, not $-26$ and $28$. Item 10: Morsella 2(c), $f_{zz}$ at $(-1, 0, -\frac{4}{3})$ is $-4$, so the point is a maximum and not a saddle point. Item 11: Damascelli 8, the determinant is $16$, not $-32$. Item 12: Eleuteri 5.1.19(11), $(0, 4)$ is not a critical point. Each was decided by a third computation (\texttt{photos/dati\_esercizi/analisi\_estremi\_vincolati.md}, section 6, and \texttt{analisi\_funzioni\_due\_variabili.md}, section 6).

\item \textbf{The other items.} Items 13--18: a saddle point called a minimum or a maximum. Items 19--25: a minimum called a maximum, or the other way. Items 26--28: a minimum or a maximum called a saddle point. Items 29--34: a point that is not critical. Items 35--44: the case in which the determinant is zero, passed off as a minimum, a maximum or a saddle point, or with the determinant called different from zero. Items 45--46: a determinant that decides, called zero. Items 47--57: a value or a comparison changed by little. Items 58--67: three and four variables, with the sign of one minor changed. Items 68--72: tables of numbers. Item 73: a bordered Hessian determinant with the wrong sign.

\item \textbf{Expected result.} Every statement is called false, with a counterexample: the thesis that fails and the values of its two sides. None is proved.

\item \textbf{How it was checked.} The independent oracle of \texttt{tools/fuzz/fuzz\_multivariable.py} (mode \texttt{check}) finds every statement false.

\item \textbf{How to run it.} From the top folder of the repository:
\begin{verbatim}
cc -std=c11 -O2 -Wall -Wextra -o /tmp/test_multivariable_calculus \
    tests/multivariable_calculus/test_multivariable_calculus.c && \
    /tmp/test_multivariable_calculus
\end{verbatim}
\end{enumerate}
\end{document}
